\section{Læringsutbytte}
\label{sec:laering}

Prosjektet har gitt oss et bredere teknisk fundament enn vi hadde da vi
startet. Kanskje vel så viktig har vi lært hvordan man jobber gjennom et
større prosjekt over tid, med eksterne avhengigheter, uklare krav og
problemer som ikke alltid har en åpenbar løsning.

Den mest håndfaste tekniske læringen er C\# og .NET. Ingen i gruppen hadde brukt språket eller rammeverket før prosjektstart,
og på 19 uker har vi gått fra å lese tutorials til å levere et API med
testdekning, dokumentasjon og containerisert deployment. Det inkluderer
å forstå dependency injection, middleware-pipelinen, konfigurasjonshåndtering
og hvordan man bygger ut et prosjekt på en måte som er vedlikeholdbar over
tid. Mye av denne læringen kom gjennom refaktoreringer underveis. Ting vi
skrev tidlig i prosjektet ville vi gjort annerledes i dag, og selve det å se
kvaliteten på koden vår øke etter hvert har vært nyttig i seg selv.

WebRTC var et nytt felt for oss alle. Vi gikk inn i prosjektet uten å vite
hva STUN- og TURN-servere er, hva ICE gjør, eller hvorfor en WebRTC-adapter
er nødvendig på klientsiden. Det tok tid å bygge opp en grunnleggende
forståelse, og mye av læringen kom fra feilsøking i situasjoner der
ingenting fungerte og vi ikke umiddelbart visste hvorfor. Det var
frustrerende mens det sto på, men gir oss nå en konkret forståelse av en
teknologi som er sentral i moderne sanntidskommunikasjon.

På infrastruktursiden har vi lært praktisk container-orkestrering med Docker
Compose. Det handler ikke bare om hvordan man skriver en
\texttt{compose.yaml}, men også om hvordan health checks faktisk oppfører
seg, hva som skjer når en container markeres unhealthy, og hvordan tjenester
finner hverandre over interne nettverk. Observability-stacken med Prometheus
og Grafana ga oss et nytt perspektiv. Det er én ting å skrive kode som
fungerer, og noe ganske annet å ha innsyn i hvordan koden oppfører seg når
den faktisk kjører. Å bygge dashboards som viser noe nyttig krever at man
først har bestemt seg for hva man trenger å vite, og det er en disiplin i
seg selv.

Vi har også fått praktisk erfaring med autentisering og autorisasjon på et
mer reelt nivå enn i tidligere skoleprosjekter. Å sette opp en
Keycloak-instans fra bunnen av, konfigurere realms og klienter, og koble den
sammen med JWT-validering i API-et var en større læringskurve enn vi
forventet. Vi forstår nå mye bedre hva som faktisk skjer når en token
presenteres på et endepunkt. Hva som valideres, hva som ikke valideres, og
hvilke konfigurasjonsvalg som er sikre nok i ulike situasjoner.

Den minst tekniske men kanskje viktigste læringen handler om arbeidsprosess.
Vi har lært å jobbe mer profesjonelt ved å bryte arbeid ned i oppgaver,
bruke pull requests og kodegjennomgang, holde branches ryddige, og skrive
commit-meldinger som gir mening senere. Vi har jobbet mot en ekstern
oppdragsgiver med asynkron kommunikasjon, noe som tvang oss til å formulere
spørsmål skriftlig og presist i stedet for å avklare ting underveis.
Når vi var avhengige av tredjepartsleverandører som svarte tregt, måtte vi
også lære oss å parkere problemer og jobbe videre på andre fronter i stedet
for å bli sittende fast.
